\section{The bit test and its smoothed correlation}
\label{sec:inner-smoothing}

We first define the inner test, including the distribution of its queries.
Its soundness statement is uniform in the size of the affine alphabet.
This uniformity will permit us to select the inner parameters before the
outer repetition parameters.

\subsection{Folded proofs, face arrays, and queries}

For a nonempty affine alphabet $A$, a \emph{folded bit proof} is a function
\[
 f:\{P:A\to\{1,-1\}\}\longrightarrow\{1,-1\},
 \qquad f(-P)=-f(P).
\]
A query is the Boolean function $P$ itself: two descriptions that give the
same function on $A$ must receive the same answer. The intended proof at
$a\in A$ is evaluation, $f(P)=P(a)$. The evaluation encoding is the
long code of the label, and the displayed antisymmetry is the standard
folding relation~\cite{BGS1998,Hastad2001}. Nonemptiness of $A$ ensures that
$P\ne-P$, so the folding requirement is consistent.

We allow a finite random index $\xi$ with arbitrary probability weights to
select an alphabet $A_\xi$ and a folded proof $f_\xi$. Conditional on $\xi$,
all sampling below is performed on its indicated alphabet. Expectations
include $\xi$ unless a conditioning is stated. Thus a bound on the average
test failure will not be replaced by a bound for every individual index.
We usually suppress $\xi$ and write $A,f$.

Let $m,n\ge2$ be integers, to be selected later. Write
\[
 X=[n]^m,\qquad N=n^m,\qquad
 \mathcal R=\{(i,y):i\in[m],\ y\in[n]^{[m]\setminus\{i\}}\},
 \qquad M=|\mathcal R|=mn^{m-1}.
\]
Here $[n]^S$ denotes coordinate tuples indexed by $S$. If $x\in X$, then
$x_{-i}$ is the tuple obtained by deleting coordinate $i$, and $x_{-ij}$
deletes both $i$ and $j$. A \emph{row} is one entry of an array
$(B^i_y)_{(i,y)\in\mathcal R}$. We take these entries independently and
uniformly from $D(A)$ and put
\begin{equation}\label{eq:face-array}
 B(x)=\sum_{i=1}^m B^i_{x_{-i}}\in D(A),
 \qquad B(a)=(B(x)(a))_{x\in X}.
\end{equation}
All sums of affine functions or binary arrays are over $\F_2$. The
\emph{face code} is the subspace
\begin{equation}\label{eq:face-code}
 \C=\left\{z\in\F_2^X:
 z_x=\sum_{i=1}^m z^i_{x_{-i}}
 \text{ for some }z^i\in\F_2^{[n]^{[m]\setminus\{i\}}}\right\}.
\end{equation}
In particular, $B(a)\in\C$ for every $a\in A$. The code depends only on
$m,n$, even when the index $\xi$ changes the alphabet.

For $u\in\R^X$ and $d\in\R^{\C}$ define a table on $\C$ by
\begin{equation}\label{eq:inner-query}
 P_{u,d}(z)=\sgn\left(\frac1{\sqrt N}
                  \sum_{x\in X}u_x(-1)^{z_x}+d_z\right).
\end{equation}
Use a fixed convention at zero. Its pullback is the Boolean query
$P_{u,d}(B):a\mapsto P_{u,d}(B(a))$. If $P$ is a table on $\C$ and
$z\in\C$, write $P^{z,+}$ and $P^{z,-}$ for the same table with its entry
at $z$ forced to $1$ and $-1$, respectively. Repeated codewords among the
labels present no difficulty: this operation changes the query value at
every label with that codeword.

Let $0<\eta\le\sigma<1$ and $0<\delta<1$. Draw independent standard
Gaussian vectors $c,g\in\R^X$ and $\lambda\in\R^{\C}$, independently of
$\xi,B$, and set
\[
 P=P_{c+\sigma g,\eta\lambda}.
\]
The following tests fail precisely when the two indicated proof answers
are unequal. A denominator is a positive weight normalization; it need
not itself be a failure probability.
\begin{enumerate}[label=\textup{(T\arabic*)},leftmargin=*]
\item Compare $f(P_{c,0}(B))$ and $f(P(B))$, with denominator
      $b_1=10\sigma$.
\item Independently for each row, replace $B^i_y$ by a fresh uniform
      element of $D(A)$ with probability $\delta^2/m$, obtaining $B'$.
      Compare $f(P(B))$ and $f(P(B'))$, using the same $c,g,\lambda$ and
      the same table $P$ on both sides. The denominator is
      $b_2=10(\delta+\eta)$.
\item Draw $z$ uniformly from $\C$, independently of the other random
      choices, and compare $f(P^{z,+}(B))$ and $f(P^{z,-}(B))$.
      The denominator is $b_3=|\C|^{-1}$.
\end{enumerate}
Write $\eps_j$ for the failure probability of test \textup{(T$j$)}.
Although the tables are defined at ties, every unforced threshold here
has a continuous Gaussian input of positive variance, so the choice at
ties does not affect any of these probabilities.

\begin{lemma}[Honest proof budgets]\label{lem:honest}
For any choice of labels $a_\xi\in A_\xi$, the evaluation proofs
$f_\xi(P)=P(a_\xi)$ satisfy $\eps_j\le b_j$ for $j=1,2,3$.
In fact,
\[
 \eps_1\le 2\sigma,\qquad
 \eps_2\le2\delta+2\eta,\qquad
 \eps_3=|\C|^{-1}.
\]
\end{lemma}
\begin{proof}
Condition on the index and its label $a$. For any $z\in\C$, put
\[
 v_z=N^{-1/2}\bigl((-1)^{z_x}\bigr)_{x\in X};
 \qquad \norm{v_z}_{\R^X}=1.
\]
For the first test, conditional on $B$, the input
$\ip{c}{v_{B(a)}}$ is standard Gaussian, and its added perturbation
$\sigma\ip{g}{v_{B(a)}}+\eta\lambda_{B(a)}$ is an independent Gaussian
of standard deviation $\sqrt{\sigma^2+\eta^2}$. The sign identity
\eqref{eq:gaussian-signs} gives disagreement probability
\[
 \frac1\pi\arctan\sqrt{\sigma^2+\eta^2}
 \le\sqrt{\sigma^2+\eta^2}\le2\sigma.
\]

For the second test, let $D$ be the relative Hamming distance between
$B(a)$ and $B'(a)$. For each $x$, the evaluation $B(x)(a)$ can change only
if at least one of its $m$ rows was selected for resampling. A union bound
therefore gives $\E D\le\delta^2$. Consequently
\[
 \E\norm{v_{B(a)}-v_{B'(a)}}_{\R^X}
 =2\E\sqrt D\le2\delta.
\]
Conditional on $B,B'$, apply \eqref{eq:gaussian-signs} to the isotropic
Gaussian $c+\sigma g$, initially omitting the table perturbations. Its
two signs disagree with probability at most
$\norm{v_{B(a)}-v_{B'(a)}}$. Restoring each perturbation
$\eta\lambda_{B(a)}$ or $\eta\lambda_{B'(a)}$ changes the corresponding
sign with probability at most
$\eta/\sqrt{1+\sigma^2}\le\eta$. The union bound gives
$\eps_2\le2\delta+2\eta$. Independence of the two table perturbations
is not needed, so this reasoning also covers $B(a)=B'(a)$.

In the third test, the two evaluation answers differ if and only if
$z=B(a)$. Uniform sampling of $z$ gives exactly $|\C|^{-1}$. All bounds
hold conditional on the chosen label and can be averaged over $\xi$.
\end{proof}

\subsection{The decoding statement}

For every folded proof define
\begin{equation}\label{eq:smoothed-gradient}
 F_B(c)=\E_{g,\lambda} f(P_{c+\sigma g,\eta\lambda}(B)),
 \qquad h_x(B,c)=\sqrt N\,\partial_{c_x}F_B(c).
\end{equation}
These derivatives exist everywhere by Gaussian convolution, as also
verified below. For real arrays on $X$ our conventions are
\[
 \ip{v}{w}_x=\frac1N\sum_{x\in X}v_xw_x,
 \qquad \norm{v}_{2,x}^2=\frac1N\sum_{x\in X}|v_x|^2.
\]
If a row $(i,x_{-i})$ is removed, the remaining rows are its
\emph{background}. For any background write
$B[B^i_{x_{-i}}\leftarrow b]$ for the array with that row filled by
$b\in D(A)$.

\begin{theorem}[Inner decoding]\label{thm:inner}
For every rational $J\ge1$, there are effectively selectable integers
$m,n\ge2$ and positive rational constants $\sigma,\delta,\eta,p,\theta<1$,
with $\eta\le\sigma$, depending only on $J$, with the following property.
For any finite weighted family of nonempty affine alphabets and folded
bit proofs, suppose
\[
 \eps_j\le Jb_j\qquad(j=1,2,3).
\]
Draw $\xi$, a standard Gaussian $c\in\R^X$, independent uniform
$x\in X$ and $i\in[m]$, and all background rows independently and
uniformly in $D(A_\xi)$. With probability at least $p$, the function
\[
 b\longmapsto h_x(B[B^i_{x_{-i}}\leftarrow b],c)
\]
has a Fourier coefficient of absolute value at least $\theta$ at an odd
frequency of $D(A_\xi)$. Equivalently, the event is
\begin{equation}\label{eq:inner-atom-event}
 \max_{\alpha\in D(A_\xi)^*: \alpha(1)=1}
 \left|\E_{b\in D(A_\xi)}
 h_x(B[B^i_{x_{-i}}\leftarrow b],c)(-1)^{\alpha(b)}\right|
 \ge\theta.
\end{equation}
For every folded proof, without any test assumption, one has pointwise
\begin{equation}\label{eq:gradient-bounds}
 \norm{h(B,c)}_{2,x}\le\sigma^{-1},\qquad
 |h_x(B,c)|\le\sqrt N/\sigma.
\end{equation}
\end{theorem}

The rest of this Section proves a low-degree correlation estimate.
Section~\ref{sec:inner-box} turns that estimate into
\eqref{eq:inner-atom-event}; the parameter choices in
Section~\ref{sec:inner-completion} complete the proof of
Theorem~\ref{thm:inner}.

\subsection{Quantitative smoothing and degree projections}

For the rest of this Section fix the following parameters, before choosing
$m,n$:
\begin{equation}\label{eq:smoothing-parameters}
 \sigma=\frac1{2560J},\qquad
 \delta=\eta=\frac{\sigma^2}{2560J},\qquad
 R=\lceil\delta^{-2}\rceil,\qquad
 T=\left\lceil\frac{32J}{\eta\sigma^4}\right\rceil.
\end{equation}
Choose an integer $s\ge1$ such that
\begin{equation}\label{eq:hermite-cutoff-choice}
 s+2>\sigma^{-2},\qquad
 (s+2)(1+\sigma^2)^{-(s+2)}\le\frac1{16}.
\end{equation}
Such an $s$ is found by increasing integers and testing rational
inequalities. The sequence $k(1+\sigma^2)^{-k}$ tends to zero, and is
decreasing once $k>\sigma^{-2}$, since the ratio of its next term to its
current term is $(1+1/k)/(1+\sigma^2)<1$. Thus these choices are effective
and imply
\begin{equation}\label{eq:hermite-tail-choice}
 \sup_{k\ge s+2} k(1+\sigma^2)^{-k}\le\frac1{16}.
\end{equation}
All constants in \eqref{eq:smoothing-parameters}--\eqref{eq:hermite-cutoff-choice}
are independent of the alphabet, its proof, and the integers $m,n$.

The Fourier degree of a function of $B$ is its degree in the independent
\emph{rows}. More precisely, a product character is
\[
 \chi_{\boldsymbol\alpha}(B)
 =\prod_{(i,y)\in\mathcal R}(-1)^{\alpha^i_y(B^i_y)},
 \qquad \alpha^i_y\in D(A)^*,
\]
and its degree is the number of nonzero $\alpha^i_y$. This counts nonzero
rows, regardless of the dimension of each row's vector space. Gaussian
degree means total degree $|I|$ in the product Hermite basis $\psi_I(c)$.
Conditional on each $\xi$, let $\Pi_B$ project onto Fourier row degree at
most $Rm$, and let $\Pi_c$ project onto Hermite degree at most $s$.
They act on different variables and commute. Define
\begin{equation}\label{eq:low-degree-field}
 L=\Pi_B\Pi_c h.
\end{equation}
The definition is coordinatewise in $x$. Both projections also act
orthogonally in the Hilbert space averaged over the arbitrary weights
of $\xi$.

\begin{lemma}[Gradient energy and the two cutoffs]\label{lem:smoothing-cutoffs}
Assume $\eps_1\le Jb_1$ and $\eps_2\le Jb_2$, and use
\eqref{eq:smoothing-parameters}--\eqref{eq:hermite-cutoff-choice}.
For every $m,n\ge2$,
\begin{align}
 \E\norm{h}_{2,x}^2&\ge\frac34,\label{eq:gradient-lower}\\
 \E\norm{(I-\Pi_c)h}_{2,x}^2&\le\frac1{16},\label{eq:c-tail}\\
 \E\norm{(I-\Pi_B)h}_{2,x}^2&\le\frac1{32},\label{eq:B-tail}\\
 \E\ip{L}{h}_x=\E\norm{L}_{2,x}^2&\ge\frac12,
 \qquad \E\norm{L}_{2,x}^2\le\sigma^{-2}.
 \label{eq:L-energy}
\end{align}
\end{lemma}
\begin{proof}
Set
$Q_B(u)=\E_\lambda f(P_{u,\eta\lambda}(B))$.
This is a bounded measurable function of $u$: it is obtained from finitely
many threshold signs followed by a function on a finite set. Gaussian
convolution and \eqref{eq:gaussian-derivative} give the identity
\begin{equation}\label{eq:h-convolution}
 h_x(B,c)=\frac{\sqrt N}{\sigma}
 \E_{g,\lambda}g_xf(P_{c+\sigma g,\eta\lambda}(B)).
\end{equation}
For fixed $\xi,B,c$, the functions $g_x$ form an orthonormal family in
the joint probability space of $(g,\lambda)$. Bessel's inequality gives
\[
 \norm{h}_{2,x}^2
 =\frac1{\sigma^2}\sum_x
 \left|\E_{g,\lambda}g_xf(P(B))\right|^2
 \le\frac1{\sigma^2}\E_{g,\lambda}|f(P(B))|^2
 =\sigma^{-2}.
\]
The individual-coordinate bound in \eqref{eq:gradient-bounds} follows
as well. In particular all gradients used here are square-integrable.

Negating $u$ and changing $\lambda$ to $-\lambda$ negates every table
entry almost surely. Indeed, conditional on $u$, each entry has an
independent continuous perturbation of variance $\eta^2>0$. Folding
therefore implies $Q_B(-u)=-Q_B(u)$. Symmetry of $g$ gives
$F_B(-c)=-F_B(c)$, so $\E_c F_B(c)=0$ for every $\xi,B$.
Let $f_0(B,c)=f(P_{c,0}(B))$, which has absolute value one. Jensen's
inequality and the first test give
\[
 \E|F_B(c)-f_0(B,c)|^2
 \le\E_{\xi,B,c,g,\lambda}|f(P(B))-f_0(B,c)|^2
 =4\eps_1\le40J\sigma\le\frac1{64}.
\]
The triangle inequality in the joint $L^2$ space consequently gives
$\norm{F}_2\ge1-1/8=7/8$ and $\E F^2\ge49/64>3/4$.
By the mean-zero Gaussian energy identity \eqref{eq:hermite-energy},
\[
 \E\norm{h}_{2,x}^2=\E\sum_x|\partial_{c_x}F_B(c)|^2
 \ge\E|F_B(c)|^2\ge\frac34,
\]
proving \eqref{eq:gradient-lower}. This argument only uses the averaged
failure bound.

For the Gaussian cutoff put
\[
 \rho=(1+\sigma^2)^{-1/2},\qquad
 G_B(v)=Q_B(\sqrt{1+\sigma^2}\,v).
\]
Then $F_B=T_\rho G_B$ and $|G_B|\le1$. If
$G_B=\sum_I a_I(B)\psi_I$, its noise average has coefficients
$\rho^{|I|}a_I(B)$. For precision about the derivative, Gaussian
integration by parts and the Hermite recurrence give
\[
 \int\partial_{c_x}F_B(c)\psi_J(c)\dd\gamma_N(c)
 =\sqrt{J_x+1}\int F_B(c)\psi_{J+e_x}(c)\dd\gamma_N(c).
\]
Integration by parts is legitimate because $F_B$ and its first
derivatives are bounded, and a Gaussian times any polynomial has
integrable tails. This identifies the actual derivative coefficients
with those obtained by formally differentiating the Hermite expansion.
Parseval and $\sum_I|a_I(B)|^2\le1$ now yield, for every $\xi,B$,
\[
 \E_c\norm{(I-\Pi_c)h}_{2,x}^2
 =\sum_{|I|\ge s+2}|I|\rho^{2|I|}|a_I(B)|^2
 \le\sup_{k\ge s+2} k(1+\sigma^2)^{-k}\le\frac1{16}.
\]
The shift from $s+1$ to $s+2$ occurs because differentiating a degree
$k$ Hermite polynomial gives degree $k-1$. This proves \eqref{eq:c-tail}
uniformly in $N$.

For the row cutoff, apply Bessel's inequality to the difference of the
two representations \eqref{eq:h-convolution}, with the same $g,\lambda$:
\begin{align}
 \E\norm{h(B,c)-h(B',c)}_{2,x}^2
 &\le\sigma^{-2}\E|f(P(B))-f(P(B'))|^2\notag\\
 &=4\eps_2/\sigma^2\le4Jb_2/\sigma^2=\frac1{32}.
 \label{eq:gradient-row-comparison}
\end{align}
Write $a=\delta^2/m$. Conditional expectation over the row-resampling
channel multiplies a constant character on one row by $1$, and a
nonconstant character by $1-a$: on resampling, its uniform average is
zero. Independence of rows therefore gives eigenvalue $(1-a)^k$ for
row degree $k$. Stationarity of the channel and Parseval show that a
degree-$k$ squared Fourier norm contributes
\[
 2\bigl(1-(1-a)^k\bigr)
\]
times that norm to the left side of
\eqref{eq:gradient-row-comparison}. If $k>Rm$, then
$ka>R\delta^2\ge1$. The elementary inequality
$(1-a)^k\le(1+ka)^{-1}$ follows by applying the binomial inequality to
$(1-a)^{-k}$, using $(1-a)^{-1}\ge1+a$. Hence $(1-a)^k<1/2$, and the
displayed multiplier exceeds one. The total squared Fourier norm of
row degrees $k>Rm$ is at most $1/32$, proving \eqref{eq:B-tail}.

Finally, commuting orthogonal projections give
$\E\ip{L}{h}_x=\E\norm{L}_{2,x}^2$. Their discarded parts satisfy
\[
 \E\norm{h-L}_{2,x}^2
 \le\E\norm{(I-\Pi_B)h}_{2,x}^2
     +\E\norm{(I-\Pi_c)h}_{2,x}^2
 \le\frac1{32}+\frac1{16}.
\]
Combining this with \eqref{eq:gradient-lower} gives
$\E\norm{L}_{2,x}^2\ge21/32>1/2$. Orthogonal projection and
\eqref{eq:gradient-bounds} give the upper bound in
\eqref{eq:L-energy}.
\end{proof}

\subsection{Truth-table differentiation and a product-of-faces correlation}

For a real array $v$ on $X$, define
\begin{equation}\label{eq:face-correlation}
 \mathfrak d_X(v)
 =\max_{z\in\C}\left|\frac1N\sum_{x\in X}v_x(-1)^{z_x}\right|
 =\max_{z^1,\ldots,z^m}
 \left|\E_{x\in X}v_x\prod_{i=1}^m(-1)^{z^i_{x_{-i}}}\right|.
\end{equation}
The second maximum is over the binary face arrays appearing in
\eqref{eq:face-code}; equality of the maxima follows from that definition.
In particular $\mathfrak d_X(v)\le\norm{v}_{2,x}$.

\begin{proposition}[Low-degree correlation]\label{prop:correlation}
Fix $J$ and the constants
\eqref{eq:smoothing-parameters}--\eqref{eq:hermite-cutoff-choice}.
For every $m,n\ge2$ and every finite weighted family of folded proofs
satisfying $\eps_j\le Jb_j$ for $j=1,2,3$, the field $L$ from
\eqref{eq:low-degree-field} has Fourier row degree at most $Rm$ and
Gaussian degree at most $s$, satisfies \eqref{eq:L-energy}, and obeys
\begin{equation}\label{eq:correlation-lower}
 \E\mathfrak d_X(L)\ge\frac1{4T}.
\end{equation}
All constants preceding $m,n$ are uniform in their values and in the
affine alphabets.
\end{proposition}
\begin{proof}
It remains to prove \eqref{eq:correlation-lower}. For fixed $B$ define
the bounded function on $\R^{\C}$
\[
 \Phi_B(t)=f\bigl(a\mapsto\sgn(t_{B(a)})\bigr),\qquad
 \Psi_B(t)=\E_\lambda\Phi_B(t+\eta\lambda).
\]
Gaussian differentiation in the independent truth-table coordinates
gives
\[
 \partial_{t_z}\Psi_B(t)
 =\eta^{-1}\E_\lambda\lambda_z\Phi_B(t+\eta\lambda).
\]
In \eqref{eq:smoothed-gradient}, the input to this function is
$t_z=N^{-1/2}\sum_x(c_x+\sigma g_x)(-1)^{z_x}$. The chain rule therefore
gives the exact representation
\begin{equation}\label{eq:truth-table-representation}
 h_x(B,c)=\sum_{z\in\C}\mu_z(B,c)(-1)^{z_x},\qquad
 \mu_z(B,c)=\eta^{-1}\E_{g,\lambda}\lambda_zf(P(B)).
\end{equation}
There are finitely many table coordinates, and
$|\partial_{t_z}\Psi_B|\le\eta^{-1}\E|\lambda_z|$ uniformly in $t$,
so differentiation may pass through the expectation over $g$.

We bound the absolute coefficients by the third bit test. Condition on
$\xi,B,c,g$ and on all coordinates of $\lambda$ other than $\lambda_z$.
The two numbers
\[
 f_+=f(P^{z,+}(B)),\qquad f_-=f(P^{z,-}(B))
\]
do not depend on $\lambda_z$. If
$t_z=N^{-1/2}\sum_x(c_x+\sigma g_x)(-1)^{z_x}$ and
$\varphi(u)=(2\pi)^{-1/2}e^{-u^2/2}$, elementary integration of the
one-dimensional Gaussian density yields
\begin{align*}
 \E_{\lambda_z}\lambda_zf(P(B))
 &=f_+\int_{-t_z/\eta}^{\infty}u\varphi(u)\dd u
   +f_-\int_{-\infty}^{-t_z/\eta}u\varphi(u)\dd u\\
 &=(f_+-f_-)\varphi(t_z/\eta).
\end{align*}
Its absolute value is at most
$2\mathbf1_{\{f_+\ne f_-\}}$. This calculation depends only on the
two possible table values at $z$; it does not require $B:A\to\C$ to
be injective. Taking absolute values after the other averages and
summing over $z$ gives, with $W=\sum_{z\in\C}|\mu_z|$,
\begin{equation}\label{eq:truth-table-mass}
 \E W\le\frac2\eta\sum_{z\in\C}
       \Prob[f(P^{z,+}(B))\ne f(P^{z,-}(B))]
 =\frac{2|\C|}\eta\eps_3\le\frac{2J}\eta.
\end{equation}
Thus \eqref{eq:truth-table-representation} places $h$ in $W$ times
the convex hull of the signed codeword arrays
$\{\pm((-1)^{z_x})_{x\in X}:z\in\C\}$ (with $h=0$ if $W=0$).
For every realization it follows that
\begin{equation}\label{eq:convex-hull-correlation}
 |\ip{L}{h}_x|\le W\mathfrak d_X(L).
\end{equation}

The large-$W$ contribution can be controlled without a pointwise bound
on $L$. By \eqref{eq:gradient-bounds}, Cauchy--Schwarz, Markov's
inequality, and \eqref{eq:L-energy},
\begin{align*}
 \E\bigl[\mathbf1_{\{W>T\}}|\ip{L}{h}_x|\bigr]
 &\le\sigma^{-1}\E\bigl[\mathbf1_{\{W>T\}}\norm{L}_{2,x}\bigr]\\
 &\le\sigma^{-1}(\E\norm{L}_{2,x}^2)^{1/2}
                 \Prob[W>T]^{1/2}\\
 &\le\sigma^{-2}\sqrt{\frac{2J}{\eta T}}\le\frac14.
\end{align*}
Since $\E\ip{L}{h}_x\ge1/2$, we obtain
\[
 \frac14
 \le\E\bigl[\mathbf1_{\{W\le T\}}\ip{L}{h}_x\bigr]
 \le T\E\mathfrak d_X(L),
\]
which is \eqref{eq:correlation-lower}. The low-degree and energy
claims have already been proved in Lemma~\ref{lem:smoothing-cutoffs}.
\end{proof}

At this stage $m,n$ remain free. The next Section chooses $m$ from
$R,s$, then a localization scale and $n$, and finally $p,\theta$.
It uses \eqref{eq:correlation-lower} to prove the asserted one-row
Fourier atom. In particular, the possibly large truth-table truncation
constant $T$ has been chosen before any of these later parameters.
